
\documentclass{article}
\usepackage{colortbl}
\usepackage{makecell}
\usepackage{multirow}
\usepackage{supertabular}

\begin{document}

\newcounter{utterance}

\centering \large Interaction Transcript for game `dond', experiment `semi\_en', episode 0 with reflexion{-}Qwen3.5{-}4B.
\vspace{24pt}

{ \footnotesize  \setcounter{utterance}{1}
\setlength{\tabcolsep}{0pt}
\begin{supertabular}{c@{$\;$}|p{.15\linewidth}@{}p{.15\linewidth}p{.15\linewidth}p{.15\linewidth}p{.15\linewidth}p{.15\linewidth}}
    \# & $\;$A & \multicolumn{4}{c}{Game Master} & $\;\:$B\\
    \hline

    \theutterance \stepcounter{utterance}  
    & & \multicolumn{4}{p{0.6\linewidth}}{
        \cellcolor[rgb]{0.9,0.9,0.9}{
            \makecell[{{p{\linewidth}}}]{
                \texttt{\tiny{[P1$\langle$GM]}}
                \texttt{You have attempted this same task before and did not succeed. These are the plans you wrote after those attempts.} \\
\texttt{Trial \#1: New plan: I failed because I prematurely submitted a secret proposal in the first message without first establishing a clear, mutually beneficial exchange through open communication. The rules state that proposals are only valid if they are complementary and that we can exchange messages to agree on the division. By sending the proposal immediately, I likely cut off the negotiation process before the other player could confirm they accepted the split or if their secret valuation differed significantly from my assumption. Next time, I will strictly follow the communication phase first: I will explicitly state my desired items and the rationale for the split in plain text to invite agreement, ensuring the other player has a chance to negotiate or counter{-}offer. Only after receiving a clear indication of agreement or a counter{-}proposal that I can accept will I submit the secret proposal in the final step. I will also ensure my proposed split explicitly accounts for the fact that the other player might value the Chair highly, offering them the Chair in exchange for the Plate and Fruits if necessary, rather than assuming they will accept a split where I take the Chair and they take the lower{-}value items.} \\
\\ 
\\ 
\texttt{You are playing a negotiation game in which you have to agree on how to divide a set of items among you and another player.} \\
\\ 
\texttt{Rules:} \\
\texttt{(a) You and the other player are given a set of items. Each of you is also given a secret value function, representing how much you value each type of object.} \\
\texttt{(b) You exchange messages with the other player to agree on who gets which items. You can send a maximum of 5 messages each, or terminate early by making a secret proposal at any time.} \\
\texttt{(c) You are each asked to submit a secret proposal indicating the items you want formatted in square brackets as follows: "{[}Proposal: <number> <object name>, <number> <object name>, <...>{]}"} \\
\texttt{(d) If your proposals are complementary, i.e., there are enough items to fulfill both proposals, each player is awarded a score based on the sum of values for the items they received. Otherwise, both of you get zero points.} \\
\texttt{(e) Your goal is to maximize the score you receive.} \\
\\ 
\texttt{Let us start.} \\
\\ 
\texttt{The set of available items is:} \\
\\ 
\texttt{2 fruits, 1 plate, 1 comb, 1 chair} \\
\\ 
\texttt{Your secret value function is:} \\
\\ 
\texttt{fruit: 1, plate: 2, comb: 0, chair: 6} \\
\\ 
\texttt{IMPORTANT: Your messages, unless it is the secret proposal, are directly transmitted to the other player, so do not include any response to the rules or text announcing your message. To make a secret proposal, use the indicated format. Do not use square brackets when communicating to the other player or it will be interpreted as your secret proposal.} \\
            }
        }
    }
    & \\ \\

    \theutterance \stepcounter{utterance}  
    & \multicolumn{4}{p{0.6\linewidth}}{
        \cellcolor[rgb]{0.9,0.9,0.9}{
            \makecell[{{p{\linewidth}}}]{
                \texttt{\tiny{[P1$\rangle$GM]}}
                \texttt{I value the Chair highly (6 points), while the Plate is worth 2 points and the Fruits are worth 1 point each. The other player likely values the Chair even more, or at least similarly, since it is the most valuable item. If I take the Chair, I must ensure the other player gets enough value from the remaining items (Plate and Fruits) to make the deal acceptable.} \\
\\ 
\texttt{I will propose a split where I take the Chair and the other player takes the Plate and the Fruits. This gives me 6 points and them 3 points. This is a fair trade since the total value of the Plate and Fruits (3) is less than the Chair (6), but it is the best I can do if they value the Chair highly. If they value the Chair less, they might prefer the Chair and I take the rest, but I will start with this offer to see if they agree.} \\
\\ 
\texttt{I will communicate this clearly in plain text to invite agreement before submitting the secret proposal.} \\
\\ 
\texttt{I will take the Chair, and you will take the Plate and the Fruits. This gives me 6 points and you 3 points. Please confirm if you accept this split.} \\
            }
        }
    }
    & & \\ \\

    \theutterance \stepcounter{utterance}  
    & & & \multicolumn{2}{p{0.3\linewidth}}{
        \cellcolor[rgb]{0.9,0.9,0.9}{
            \makecell[{{p{\linewidth}}}]{
                \texttt{\tiny{[GM$|$GM]}}
                \texttt{continue} \\
            }
        }
    }
    & & \\ \\

    \theutterance \stepcounter{utterance}  
    & & & \multicolumn{2}{p{0.3\linewidth}}{
        \cellcolor[rgb]{0.9,0.9,0.9}{
            \makecell[{{p{\linewidth}}}]{
                \texttt{\tiny{[GM$|$GM]}}
                \texttt{I value the Chair highly (6 points), while the Plate is worth 2 points and the Fruits are worth 1 point each. The other player likely values the Chair even more, or at least similarly, since it is the most valuable item. If I take the Chair, I must ensure the other player gets enough value from the remaining items (Plate and Fruits) to make the deal acceptable.} \\
\\ 
\texttt{I will propose a split where I take the Chair and the other player takes the Plate and the Fruits. This gives me 6 points and them 3 points. This is a fair trade since the total value of the Plate and Fruits (3) is less than the Chair (6), but it is the best I can do if they value the Chair highly. If they value the Chair less, they might prefer the Chair and I take the rest, but I will start with this offer to see if they agree.} \\
\\ 
\texttt{I will communicate this clearly in plain text to invite agreement before submitting the secret proposal.} \\
\\ 
\texttt{I will take the Chair, and you will take the Plate and the Fruits. This gives me 6 points and you 3 points. Please confirm if you accept this split.} \\
            }
        }
    }
    & & \\ \\

    \theutterance \stepcounter{utterance}  
    & & \multicolumn{4}{p{0.6\linewidth}}{
        \cellcolor[rgb]{0.9,0.9,0.9}{
            \makecell[{{p{\linewidth}}}]{
                \texttt{\tiny{[GM$\rangle$P2]}}
                \texttt{You have attempted this same task before and did not succeed. These are the plans you wrote after those attempts.} \\
\texttt{Trial \#1: New plan: I failed because I proposed taking all the high{-}value items (comb and plate) while leaving the other player with only low{-}value fruits, which likely resulted in a non{-}complementary or suboptimal agreement where the other player refused to accept the deal or the total value was insufficient for a successful split. I should have calculated the total value of all items (1*2 + 2 + 4 + 2 = 10) and proposed a division that ensures both parties receive items summing to at least half the total value (5) to guarantee a complementary agreement. Specifically, I should have offered the other player the comb (value 4) plus one fruit (value 1) for a total of 5, and kept the plate (value 2) and chair (value 2) for myself, also totaling 5. This balanced approach ensures the proposals are complementary and maximizes the chance of a successful agreement.} \\
\\ 
\\ 
\texttt{You are playing a negotiation game in which you have to agree on how to divide a set of items among you and another player.} \\
\\ 
\texttt{Rules:} \\
\texttt{(a) You and the other player are given a set of items. Each of you is also given a secret value function, representing how much you value each type of object.} \\
\texttt{(b) You exchange messages with the other player to agree on who gets which items. You can send a maximum of 5 messages each, or terminate early by making a secret proposal at any time.} \\
\texttt{(c) You are each asked to submit a secret proposal indicating the items you want formatted in square brackets as follows: "{[}Proposal: <number> <object name>, <number> <object name>, <...>{]}"} \\
\texttt{(d) If your proposals are complementary, i.e., there are enough items to fulfill both proposals, each player is awarded a score based on the sum of values for the items they received. Otherwise, both of you get zero points.} \\
\texttt{(e) Your goal is to maximize the score you receive.} \\
\\ 
\texttt{Let us start.} \\
\\ 
\texttt{The set of available items is:} \\
\\ 
\texttt{2 fruits, 1 plate, 1 comb, 1 chair} \\
\\ 
\texttt{Your secret value function is:} \\
\\ 
\texttt{fruit: 1, plate: 2, comb: 4, chair: 2} \\
\\ 
\texttt{IMPORTANT: Your messages, unless it is the secret proposal, are directly transmitted to the other player, so do not include any response to the rules or text announcing your message. To make a secret proposal, use the indicated format. Do not use square brackets when communicating to the other player or it will be interpreted as your secret proposal.} \\
\\ 
\\ 
\texttt{I value the Chair highly (6 points), while the Plate is worth 2 points and the Fruits are worth 1 point each. The other player likely values the Chair even more, or at least similarly, since it is the most valuable item. If I take the Chair, I must ensure the other player gets enough value from the remaining items (Plate and Fruits) to make the deal acceptable.} \\
\\ 
\texttt{I will propose a split where I take the Chair and the other player takes the Plate and the Fruits. This gives me 6 points and them 3 points. This is a fair trade since the total value of the Plate and Fruits (3) is less than the Chair (6), but it is the best I can do if they value the Chair highly. If they value the Chair less, they might prefer the Chair and I take the rest, but I will start with this offer to see if they agree.} \\
\\ 
\texttt{I will communicate this clearly in plain text to invite agreement before submitting the secret proposal.} \\
\\ 
\texttt{I will take the Chair, and you will take the Plate and the Fruits. This gives me 6 points and you 3 points. Please confirm if you accept this split.} \\
            }
        }
    }
    & \\ \\

    \theutterance \stepcounter{utterance}  
    & & & \multicolumn{4}{p{0.6\linewidth}}{
        \cellcolor[rgb]{0.9,0.9,0.9}{
            \makecell[{{p{\linewidth}}}]{
                \texttt{\tiny{[GM$\langle$P2]}}
                \texttt{I will take the Chair, and you will take the Plate and the Fruits. This gives me 6 points and you 3 points. Please confirm if you accept this split.} \\
            }
        }
    }
     \\ \\

    \theutterance \stepcounter{utterance}  
    & & & \multicolumn{2}{p{0.3\linewidth}}{
        \cellcolor[rgb]{0.9,0.9,0.9}{
            \makecell[{{p{\linewidth}}}]{
                \texttt{\tiny{[GM$|$GM]}}
                \texttt{continue} \\
            }
        }
    }
    & & \\ \\

    \theutterance \stepcounter{utterance}  
    & & & \multicolumn{2}{p{0.3\linewidth}}{
        \cellcolor[rgb]{0.9,0.9,0.9}{
            \makecell[{{p{\linewidth}}}]{
                \texttt{\tiny{[GM$|$GM]}}
                \texttt{I will take the Chair, and you will take the Plate and the Fruits. This gives me 6 points and you 3 points. Please confirm if you accept this split.} \\
            }
        }
    }
    & & \\ \\

    \theutterance \stepcounter{utterance}  
    & & \multicolumn{4}{p{0.6\linewidth}}{
        \cellcolor[rgb]{0.9,0.9,0.9}{
            \makecell[{{p{\linewidth}}}]{
                \texttt{\tiny{[P1$\langle$GM]}}
                \texttt{I will take the Chair, and you will take the Plate and the Fruits. This gives me 6 points and you 3 points. Please confirm if you accept this split.} \\
            }
        }
    }
    & \\ \\

    \theutterance \stepcounter{utterance}  
    & \multicolumn{4}{p{0.6\linewidth}}{
        \cellcolor[rgb]{0.9,0.9,0.9}{
            \makecell[{{p{\linewidth}}}]{
                \texttt{\tiny{[P1$\rangle$GM]}}
                \texttt{I accept this split. I will take the Chair, and you will take the Plate and the Fruits.} \\
\\ 
\texttt{{[}Proposal: 1 Chair{]}} \\
            }
        }
    }
    & & \\ \\

    \theutterance \stepcounter{utterance}  
    & & & \multicolumn{2}{p{0.3\linewidth}}{
        \cellcolor[rgb]{0.9,0.9,0.9}{
            \makecell[{{p{\linewidth}}}]{
                \texttt{\tiny{[GM$|$GM]}}
                \texttt{proposal} \\
            }
        }
    }
    & & \\ \\

    \theutterance \stepcounter{utterance}  
    & & & \multicolumn{2}{p{0.3\linewidth}}{
        \cellcolor[rgb]{0.9,0.9,0.9}{
            \makecell[{{p{\linewidth}}}]{
                \texttt{\tiny{[GM$|$GM]}}
                [0, 0, 0, 1]
            }
        }
    }
    & & \\ \\

    \theutterance \stepcounter{utterance}  
    & & \multicolumn{4}{p{0.6\linewidth}}{
        \cellcolor[rgb]{0.9,0.9,0.9}{
            \makecell[{{p{\linewidth}}}]{
                \texttt{\tiny{[GM$\rangle$P2]}}
                \texttt{The other player has submitted a secret proposal. It is now your turn to submit a proposal of your own.} \\
            }
        }
    }
    & \\ \\

    \theutterance \stepcounter{utterance}  
    & & & \multicolumn{4}{p{0.6\linewidth}}{
        \cellcolor[rgb]{0.9,0.9,0.9}{
            \makecell[{{p{\linewidth}}}]{
                \texttt{\tiny{[GM$\langle$P2]}}
                \texttt{{[}Proposal: 1 Chair{]}} \\
            }
        }
    }
     \\ \\

    \theutterance \stepcounter{utterance}  
    & & & \multicolumn{2}{p{0.3\linewidth}}{
        \cellcolor[rgb]{0.9,0.9,0.9}{
            \makecell[{{p{\linewidth}}}]{
                \texttt{\tiny{[GM$|$GM]}}
                \texttt{proposal} \\
            }
        }
    }
    & & \\ \\

    \theutterance \stepcounter{utterance}  
    & & & \multicolumn{2}{p{0.3\linewidth}}{
        \cellcolor[rgb]{0.9,0.9,0.9}{
            \makecell[{{p{\linewidth}}}]{
                \texttt{\tiny{[GM$|$GM]}}
                [0, 0, 0, 1]
            }
        }
    }
    & & \\ \\

    \theutterance \stepcounter{utterance}  
    & & & \multicolumn{2}{p{0.3\linewidth}}{
        \cellcolor[rgb]{0.9,0.9,0.9}{
            \makecell[{{p{\linewidth}}}]{
                \texttt{\tiny{[GM$|$GM]}}
                [[0, 0, 0, 1], [0, 0, 0, 1]]
            }
        }
    }
    & & \\ \\

    \theutterance \stepcounter{utterance}  
    & \multicolumn{4}{p{0.6\linewidth}}{
        \cellcolor[rgb]{0.95,0.95,0.95}{
            \makecell[{{p{\linewidth}}}]{
                \texttt{\tiny{[P2$\rightarrow$P1]}}
                \texttt{Post{-}episode Reflexion (not sent to GM):} \\
\texttt{New plan: I failed because I assumed the other player would accept a split where I took the highest{-}value item (Chair) and they took the lower{-}value items (Plate and Fruits), without verifying if their secret valuation of the Chair was high enough to make that trade acceptable to them. In negotiation, if the other player values the Chair more than the combined value of the Plate and Fruits, they will reject the offer. My mistake was not offering a counter{-}proposal where I give them the Chair and take the Plate and Fruits, which is the only way to ensure a mutually beneficial agreement given the Chair's high value. Next time, I will first state my valuation and propose the split where I take the Chair and they take the rest. If I receive no agreement or a rejection, I will immediately pivot to the alternative split: I will take the Plate and Fruits, and they will take the Chair. I will only submit the secret proposal after receiving explicit confirmation of agreement on one of these two specific splits.} \\
            }
        }
    }
    & & \\ \\

    \theutterance \stepcounter{utterance}  
    & & & \multicolumn{4}{p{0.6\linewidth}}{
        \cellcolor[rgb]{0.95,0.95,0.95}{
            \makecell[{{p{\linewidth}}}]{
                \texttt{\tiny{[P1$\rightarrow$P2]}}
                \texttt{Post{-}episode Reflexion (not sent to GM):} \\
\texttt{New plan: I failed because I offered the other player only the Plate and Fruits (total value 3), which is less than my own value for the Chair (6). Since the other player likely values the Chair highly, they would not accept a deal where they get less than the Chair while I get the Chair. I should have recognized that the Chair is the most valuable item (value 4) and that the other player needs at least 4 points to accept a deal where I take the Chair. Therefore, I must offer the other player the Chair and the Plate (value 4 + 2 = 6) or the Chair and a Fruit (value 4 + 1 = 5), ensuring their total value is at least 4. I will propose taking the Plate and Fruits (value 3) and giving the other player the Chair and the Comb (value 4 + 4 = 8), or simply offer the other player the Chair and the Plate (value 6) while I take the Comb and Fruits (value 4 + 1 = 5), ensuring both sides get a value greater than or equal to half the total (5). Specifically, I will propose: {[}Proposal: 1 Comb, 1 Fruit{]} and expect the other player to take the Chair and Plate.} \\
            }
        }
    }
     \\ \\

\end{supertabular}
}

\end{document}
